\documentclass[11pt]{article}
\usepackage[T1]{fontenc}
\usepackage{textcomp}
\usepackage[margin=0.75in]{geometry}
\usepackage{amsmath,amssymb,amsthm}
\usepackage{array,booktabs}
\usepackage{hyperref}
\setlength{\parindent}{0pt}
\newtheorem{theorem}{Theorem}
\theoremstyle{definition}
\newtheorem{definition}{Definition}
\title{Flow Proof Artifact: Rat}
\author{Generated by \texttt{flow doc proof}}
\date{Flow Proof Book}
\begin{document}
\maketitle
Rational numbers as normalized pair quotients.\\[0.5em]
\textbf{Source.} landau — *Foundations of Analysis*\\[0.5em]
\bigskip
\noindent{\large \textbf{Definition 1.} \textit{Zero is the left additive identity for rationals} \hfill Rat \cdot addition \cdot zero is the left identity}\par
\smallskip\noindent\textit{Coordinate: Rat \cdot addition \cdot zero is the left identity}\par
\smallskip\noindent\textit{Source: landau}\par
\medskip
\noindent\textbf{Goal.} Zero is the left additive identity for rationals\par
\noindent$\displaystyle \forall q \in Rat\quad 0 + q = q$\par
\medskip
\noindent
\begin{tabular}{@{} >{\raggedright\arraybackslash}p{0.06\textwidth} >{\raggedright\arraybackslash}p{0.47\textwidth} >{\raggedleft\arraybackslash}p{0.41\textwidth} @{}}
\textbf{\#} & \textbf{Proof} & \textbf{Mathematics} \\
\midrule
\textbf{1.} & We stipulate zero is the left identity for addition on Rat — this is a definition, not a derived fact. & \\[0.45em]
\textbf{2.} & This follows directly from the definition: 0 plus q equals q. Hence proven. & $\displaystyle 0 + q = q$ \\[0.45em]
\bottomrule
\end{tabular}
\label{thm:Rat-addition-zero_is_the_left_identity}
\par\medskip\hrule
\bigskip
\noindent{\large \textbf{Definition 2.} \textit{Zero is the right additive identity for rationals} \hfill Rat \cdot addition \cdot zero is the right identity}\par
\smallskip\noindent\textit{Coordinate: Rat \cdot addition \cdot zero is the right identity}\par
\smallskip\noindent\textit{Source: landau}\par
\medskip
\noindent\textbf{Goal.} Zero is the right additive identity for rationals\par
\noindent$\displaystyle \forall q \in Rat\quad q + 0 = q$\par
\medskip
\noindent
\begin{tabular}{@{} >{\raggedright\arraybackslash}p{0.06\textwidth} >{\raggedright\arraybackslash}p{0.47\textwidth} >{\raggedleft\arraybackslash}p{0.41\textwidth} @{}}
\textbf{\#} & \textbf{Proof} & \textbf{Mathematics} \\
\midrule
\textbf{1.} & We stipulate zero is the right identity for addition on Rat — this is a definition, not a derived fact. & \\[0.45em]
\textbf{2.} & This follows directly from the definition: q plus 0 equals q. Hence proven. & $\displaystyle q + 0 = q$ \\[0.45em]
\bottomrule
\end{tabular}
\label{thm:Rat-addition-zero_is_the_right_identity}
\par\medskip\hrule
\bigskip
\noindent{\large \textbf{Definition 3.} \textit{One is the left multiplicative identity for rationals} \hfill Rat \cdot multiplication \cdot one is the left identity}\par
\smallskip\noindent\textit{Coordinate: Rat \cdot multiplication \cdot one is the left identity}\par
\smallskip\noindent\textit{Source: landau}\par
\medskip
\noindent\textbf{Goal.} One is the left multiplicative identity for rationals\par
\noindent$\displaystyle \forall q \in Rat\quad 1 \cdot q = q$\par
\medskip
\noindent
\begin{tabular}{@{} >{\raggedright\arraybackslash}p{0.06\textwidth} >{\raggedright\arraybackslash}p{0.47\textwidth} >{\raggedleft\arraybackslash}p{0.41\textwidth} @{}}
\textbf{\#} & \textbf{Proof} & \textbf{Mathematics} \\
\midrule
\textbf{1.} & We stipulate one is the left identity for multiplication on Rat — this is a definition, not a derived fact. & \\[0.45em]
\textbf{2.} & This follows directly from the definition: 1 times q equals q. Hence proven. & $\displaystyle 1 \cdot q = q$ \\[0.45em]
\bottomrule
\end{tabular}
\label{thm:Rat-multiplication-one_is_the_left_identity}
\par\medskip\hrule
\end{document}
